\chapter*{付録}
\addcontentsline{toc}{chapter}{付録}

\section*{付録A：特徴量の詳細な数式定義}
\addcontentsline{toc}{section}{付録A：特徴量の詳細な数式定義}

本研究で使用する10次元特徴ベクトルの詳細な数式定義を示す.
各特徴量は150サンプル（サンプリング周波数1000Hzで150ms相当）の時間窓から計算される.

\subsection*{時間領域特徴量}

\begin{enumerate}
    \item \textbf{平均絶対値 (MAV: Mean Absolute Value)}
    \begin{equation}
        \text{MAV} = \frac{1}{N} \sum_{i=1}^{N} |x_i|
    \end{equation}
    
    \item \textbf{分散 (Variance / Hjorth Activity)}
    \begin{equation}
        \sigma^2 = \frac{1}{N-1} \sum_{i=1}^{N} (x_i - \mu)^2
    \end{equation}
    
    \item \textbf{平均波形長 (AW: Average Waveform Length)}
    \begin{equation}
        \text{AW} = \frac{1}{N-1} \sum_{i=1}^{N-1} |x_{i+1} - x_i|
    \end{equation}
    
    \item \textbf{Hjorth Complexity}
    
    信号の複雑さを表す指標で,以下のように定義される：
    \begin{equation}
        \text{Mobility}(x) = \sqrt{\frac{\text{Var}(\frac{dx}{dt})}{\text{Var}(x)}}
    \end{equation}
    \begin{equation}
        \text{Complexity} = \frac{\text{Mobility}(\frac{dx}{dt})}{\text{Mobility}(x)}
    \end{equation}
\end{enumerate}

\subsection*{周波数領域特徴量}

FFT（高速フーリエ変換）により周波数スペクトルを計算し,以下の特徴量を抽出する.
入力信号にはHann窓を適用し,スペクトル漏れを抑制する.

\begin{enumerate}
    \setcounter{enumi}{4}
    \item \textbf{スペクトル重心周波数 (CPS: Center Power Spectrum)}
    \begin{equation}
        \text{CPS} = \frac{\sum_{k=1}^{N/2} f_k \cdot P_k}{\sum_{k=1}^{N/2} P_k}
    \end{equation}
    ここで,\(f_k\)は周波数ビン,\(P_k\)はパワースペクトル密度を表す.
    
    \item \textbf{低周波帯エネルギー比 (LFO: Low Frequency Ratio)}
    \begin{equation}
        \text{LFO} = \frac{\sum_{f \leq 150\text{Hz}} P_f}{\sum_{k=1}^{N/2} P_k}
    \end{equation}
    150Hz以下の周波数帯域のパワーが全体に占める割合を表す.
\end{enumerate}

\subsection*{自己回帰モデル係数}

\begin{enumerate}
    \setcounter{enumi}{6}
    \item \textbf{AR係数 (AR1--AR4)}
    
    4次自己回帰モデルの係数をYule-Walker法により推定する：
    \begin{equation}
        x_t = \sum_{j=1}^{4} a_j x_{t-j} + \epsilon_t
    \end{equation}
    
    自己相関関数\(r[k]\)を用いて、以下のToeplitz方程式を解く：
    \begin{equation}
        \begin{bmatrix}
        r[0] & r[1] & r[2] & r[3] \\
        r[1] & r[0] & r[1] & r[2] \\
        r[2] & r[1] & r[0] & r[1] \\
        r[3] & r[2] & r[1] & r[0]
        \end{bmatrix}
        \begin{bmatrix}
        a_1 \\ a_2 \\ a_3 \\ a_4
        \end{bmatrix}
        =
        \begin{bmatrix}
        r[1] \\ r[2] \\ r[3] \\ r[4]
        \end{bmatrix}
    \end{equation}
    
    これにより4つの係数\(a_1, a_2, a_3, a_4\)が得られる.
\end{enumerate}

\section*{付録B：ニューラルネットワークの詳細実装}
\addcontentsline{toc}{section}{付録B：ニューラルネットワークの詳細実装}

本研究で使用するニューラルネットワークの詳細な層構成と処理フローを示す.

\subsection*{処理フロー}

図\ref{fig:network_architecture}に示すように,Core MLのInnerProduct層とActivation層を組み合わせて
実装されている.ネットワークは以下の処理フローで構成される：

\begin{enumerate}
    \item \textbf{入力層（features）}：10次元の特徴ベクトル$\mathbf{x} \in \mathbb{R}^{10}$を受け取る.
    
    \item \textbf{第1のInnerProduct層}：10次元入力を24次元に変換
    \begin{equation}
        \mathbf{h}_{\text{pre}} = W_1 \mathbf{x} + \mathbf{b}_1
    \end{equation}
    ここで,$W_1 \in \mathbb{R}^{24 \times 10}$,$\mathbf{b}_1 \in \mathbb{R}^{24}$である.
    この層は\textbf{固定層}であり,デバイス上での学習時には更新されない.
    
    \item \textbf{Activation層（ReLU）}：非線形変換を適用
    \begin{equation}
        \mathbf{h} = \text{ReLU}(\mathbf{h}_{\text{pre}}) = \max(0, \mathbf{h}_{\text{pre}})
    \end{equation}
    
    \item \textbf{第2のInnerProduct層（Updatable）}：24次元を4クラスのlogitsに変換
    \begin{equation}
        \mathbf{z} = W_2 \mathbf{h} + \mathbf{b}_2
    \end{equation}
    ここで,$W_2 \in \mathbb{R}^{4 \times 24}$,$\mathbf{b}_2 \in \mathbb{R}^{4}$である.
    この層は\textbf{更新可能層（Updatable）}であり,デバイス上での継続学習の対象となる.
    
    \item \textbf{Softmax層}：logitsをクラス確率に変換
    \begin{equation}
        p_i = \frac{\exp(z_i)}{\sum_{j=1}^{4} \exp(z_j)}, \quad i = 1, 2, 3, 4
    \end{equation}
    
    \item \textbf{出力層（stringClassLabels）}：確率分布$\mathbf{p} = [p_1, p_2, p_3, p_4]^T$と
    最大確率を持つクラスラベル（classLabel）を出力する.
\end{enumerate}

\subsection*{計算量の詳細}

順伝播時の浮動小数点演算回数（FLOPs）は：
\begin{itemize}
    \item 隠れ層：$24 \times (10 \times 2) = 480$ FLOPs（積和算）
    \item 出力層：$4 \times (24 \times 2) = 192$ FLOPs
    \item \textbf{合計：約672 FLOPs}
\end{itemize}

\subsection*{学習アルゴリズム}

損失関数にはカテゴリカル交差エントロピーを使用し,最適化手法にはSGD（Stochastic Gradient Descent）
またはAdam optimizerを採用している：

\begin{equation}
    \mathcal{L} = -\sum_{i=1}^{4} y_i \log(p_i)
\end{equation}

ここで,$y_i$は真のラベルのone-hotエンコーディング,$p_i$はモデルが出力する確率である.

\subsection*{Transfer Learning的アプローチ}

隠れ層を固定し出力層のみを更新する設計には,以下の利点がある：

\begin{itemize}
    \item \textbf{学習の安定性}：隠れ層を固定することで,既に獲得された汎用的な特徴表現を保持し,
    過学習を抑制できる.
    
    \item \textbf{計算効率}：更新対象のパラメータが100個（総パラメータの27.5\%）に限定されるため,
    勾配計算とパラメータ更新の計算コストが削減される.
    
    \item \textbf{Transfer Learning的アプローチ}：隠れ層が共通の特徴抽出器として機能し,
    出力層のみを個人適応させることで,少量のデータでも効果的な学習が可能となる.
\end{itemize}

